\chapter{多无人机协同观测通信建模}
\label{ch:foundations}

本章建立多无人机协同观测信息回传的系统通信模型，为后续各章算法设计提供统一的分析基础。首先介绍动态无线链路、观测消息表示与确认反馈过程，随后明确定义三类通信资源约束以及任务按时正确完成条件，说明路由、消息选择与发送控制如何共同决定地面中心的最终接收结果，最后介绍后续章节使用的优化与决策方法基础。

\section{无人机网络与数据传输}

将多无人机通信网络建模为时变有向图 $\mathcal G(t)=(\mathcal V,\mathcal E(t))$，其中节点集 $\mathcal V$ 包含所有无人机与地面中心，边集 $\mathcal E(t)$ 表示时刻 $t$ 存在的可用无线链路。设节点 $i$ 的空间位置为 $\bm x_i(t)$，节点 $i$ 与 $j$ 之间的欧氏距离为
\begin{equation}
d_{ij}(t)=\|\bm x_i(t)-\bm x_j(t)\|_2.
\label{eq:base:distance}
\end{equation}
随着无人机飞行与无线信道变化，节点 $i$ 的一跳邻居集合 $\mathcal N_i(t)=\{j:(i,j)\in\mathcal E(t)\}$ 也会不断更新。采用香农信道容量模型，链路 $(i,j)$ 的瞬时传输速率表示为
\begin{equation}
r_{ij}(t)=W_{ij}(t)\log_2\left(1+
\frac{P_i(t)h_{ij}(t)}{I_j(t)+N_j(t)}\right),
\label{eq:base:rate}
\end{equation}
其中 $W_{ij}$、$P_i$ 和 $h_{ij}$ 分别表示信道带宽、节点发射功率及信道增益，$I_j$ 和 $N_j$ 分别为接收端受到的并发信号干扰与背景噪声。由式~\eqref{eq:base:rate} 可以看出，提高发射功率时，传输速率只按对数规律缓慢上升，而发送能耗却会线性增加，同时还会加重对周边无线链路的干扰。此外，在多跳网络中，同一份观测数据往往需要经过多次中继转发才能到达地面中心，因此本文在路由研究中以全网累计转发吞吐量为主要业务指标，并结合丢包率与网络生命周期共同评估转发效果。

每个无人机节点都配有本地数据缓存，用来暂存尚未发出的业务数据。记节点 $i$ 在时刻 $t$ 的数据队列积压为 $Q_i(t)$，新到达数据量为 $A_i(t)$，其中同时包含节点本地新采集的数据与上游邻居转发过来的数据，有效服务量为 $\mu_i(t)$，则队列长度按如下关系更新
\begin{equation}
Q_i(t+1)=\left[Q_i(t)-\mu_i(t)\right]^++A_i(t),
\qquad [x]^+=\max\{x,0\}.
\label{eq:base:queue}
\end{equation}
当某节点的队列积压持续增加时，说明当前链路的转发能力跟不上数据到达速度，因此队列长度可以直接作为路由决策中的拥塞参考信号。能耗方面，一次持续时长为 $\tau$、发射功率为 $P^{\rm tx}$ 的发送过程消耗的能量为 $P^{\rm tx}\tau$。若节点在当前阶段的总通信耗电为 $e_i(t)$，则其剩余通信电量满足
\begin{equation}
E_i(t+1)=E_i(t)-e_i(t),\qquad 0\le e_i(t)\le E_i(t).
\label{eq:base:energy}
\end{equation}
路由研究通过离散事件仿真统计各节点的实际通信耗电，语义选择研究重点统计消息的词元长度与传输代价，可靠交付研究则同时计入数据帧与确认帧的发送耗电。

无人机集群中各节点的剩余电量分布直接关系到回传路径能维持多久。如果少数靠近地面中心的关键中继耗电过快，这些节点提前退出网络就会导致路径中断；如果为了平衡电量而强行让所有节点平均分担转发，数据往往会被导向距离较远或干扰较重的链路，反而拖累全网转发速率。因此，能量均衡应以保障实际转发效率为前提，不能只用节点间的电量方差大小来评判路由好坏。

除链路速率与电量外，消息长度和确认机制也会影响数据传输。对于长度为 $B_m$ 比特的消息，在有效速率 $r_m$ 下的单次发送时长为 $B_m/r_m$。当消息较长或链路质量变差时，单次发送耗时就会明显增加。一方面，无人机在飞行过程中随时可能移出通信范围，单次发送时间过长，更容易在消息传完前发生断链丢包；另一方面，每项观测任务都有固定的截止时间，单次发送占用的时间越多，在截止前留给发送其他消息或重传丢失消息的时间就越少，任务也越容易超时。

\begin{figure}[htbp]
\centering
\begin{tikzpicture}[font=\small,>=Stealth]
\draw[->] (0,2.8) -- (12,2.8) node[right]{时间};
\draw[->] (0,0) -- (12,0);
\node[left] at (0,2.8){发送端};\node[left] at (0,0){接收端};
\fill (1,2.8) circle(1.5pt) node[above]{开始发送};
\fill (4,2.8) circle(1.5pt) node[above]{尝试结束};
\fill (4,0) circle(1.5pt) node[below,align=center]{载荷接收结果\\生成累计确认};
\draw[->,thick] (1,2.8)--node[above,sloped]{数据传输}(4,0);
\draw[->,thick,dashed] (4,0)--node[above,sloped]{反馈发送与延迟}(9.5,2.8);
\fill (9.5,2.8) circle(1.5pt) node[above]{收到确认};
\draw[decorate,decoration={brace,amplitude=4pt}] (4,3.6)--(9.5,3.6) node[midway,above=6pt,font=\footnotesize]{接收状态尚未完全确定};
\end{tikzpicture}

\bicaption{一次发送尝试中的数据事件与反馈事件}{Data and feedback events during one transmission attempt}
\label{fig:base:timeline}
\end{figure}

如图~\ref{fig:base:timeline} 所示，数据发出、地面中心接收并生成确认消息（ACK）以及确认消息返回发送端，分别发生在不同时刻。由于无线信道可能丢包，确认消息在回传途中也会产生延迟甚至丢失，发送端在发完数据后无法立刻确定消息是否送达。本文采用累计确认机制，每条确认消息都会携带地面中心在生成该确认时已收到的全部消息集合。在确认回传期间，发送端可能已经发起了新的传输，且受延迟波动影响，多条确认消息还可能出现先发后至的情况。因此，发送端必须依据确认消息对应的发送序号来核对并更新接收状态，不能直接按确认到达的先后顺序来判断。

\section{观测消息与任务判断}

针对第 $k$ 项观测任务，无人机 $i$ 根据局部观测画面 $o_{i,k}$ 与当前任务查询 $q_k$，通过机载语义提取模块生成局部证据 $e_{i,k}=\phi(o_{i,k},q_k)$，其中包含局部类别预测、置信度、文本描述以及画面观测质量等信息。在多视角排查任务中，判断目标区域是否具备清晰、无遮挡的判别依据是综合研判的基础，因此本文选用五级证据可见性判别作为具体任务，类别依次为清晰可见、大部分可见、不确定、大部分遮挡和完全遮挡，用来衡量局部画面对回答当前问题的支持程度。实验使用的图像与问题取自 VQAv2 和 COCO 公开数据集，由视觉语言模型对完整图像作出的判别结果作为弱监督参考标签，各无人机节点只能看到原始图像中的局部裁剪区域。

地面中心在收到消息后，据此估计当前任务属于各类别的后验概率，将定义在类别空间 $\mathcal Y$ 上的概率分布记为语义信念
\begin{equation}
\bm b^{\rm sem}=(b_1,\ldots,b_{|\mathcal Y|}),\quad
b_y\ge0,\quad \sum_{y\in\mathcal Y}b_y=1.
\label{eq:base:semantic-belief}
\end{equation}
其中概率最大的分量对应地面中心当前的预测类别，香农熵 $H(\bm b^{\rm sem})=-\sum_yb_y\log b_y$ 用来衡量当前判断的不确定程度。单条消息自带的置信度虽然可供筛选参考，但在严重遮挡或只看到局部画面时，模型同样可能给出高置信度的错误判断。此外，同一份局部证据可以组织成类别标签、简短陈述、摘要或完整描述等不同详略程度的消息。通常消息越长，携带的细节越丰富，但占用的带宽与发送时间也越多，它对最终判断的实际帮助取决于地面中心目前已经掌握了哪些信息。

设地面中心当前已收到的消息集合为 $S$，$F_k(S)$ 表示依据集合 $S$ 得到的任务判断质量，则新加入一条消息 $m$ 带来的质量变化可写为
\begin{equation}
\Delta F_k(m\mid S)=F_k(S\cup\{m\})-F_k(S).
\label{eq:base:marginal}
\end{equation}
如果新消息的内容与已有消息高度重叠，往往很难改善当前判断；而来自未覆盖角度的互补证据则可能直接消除歧义。由此可见，同一条候选消息的价值会随着已收集合 $S$ 的变化而改变，仅靠通信前的一次性静态排序无法反映这种先后依赖关系。为此，后文的语义选择研究在每选入一条消息后都会动态更新候选评价，而交付控制研究则在筛选后的候选集上直接评估不同接收组合的任务效用。

\begin{figure}[htbp]
\centering
\begin{tikzpicture}[
 x=1mm,y=1mm,>=Stealth,
 font=\fontsize{9bp}{11.5bp}\selectfont,
 box/.style={draw,line width=.45pt,align=center,text width=30mm,
   minimum height=13mm,inner sep=1mm},
 edge/.style={->,line width=.55pt},
 note/.style={align=center,text width=30mm,inner sep=0pt}]
\path[use as bounding box] (0,0) rectangle (148,-23);
\node[box] (a) at (16.25,-7.5) {局部图像\\与任务查询};
\node[box] (b) at (54.75,-7.5) {类别判断\\简短陈述或摘要};
\node[box] (c) at (93.25,-7.5) {按预算选择消息\\根据确认安排发送};
\node[box] (d) at (131.75,-7.5) {融合已收到的信息\\形成任务判断};
\draw[edge] (a.east) -- (b.west);
\draw[edge] (b.east) -- (c.west);
\draw[edge] (c.east) -- (d.west);
\node[note] at (54.75,-20) {保留任务所需内容};
\node[note] at (93.25,-20) {控制发送与重传};
\node[note] at (131.75,-20) {检查判断结果};
\end{tikzpicture}

\bicaption{局部观测消息的生成、传输与任务判断}{Observation message generation, transmission, and task inference}
\label{fig:base:semantic-pipeline}
\end{figure}

图~\ref{fig:base:semantic-pipeline} 展示了局部证据从生成、筛选、回传到融合判断的完整过程。发送前，系统结合预算权衡各条候选消息的价值并挑选待发内容；消息发出后尚未收到确认时，发送端需要在重传与等待之间作出决定；一旦某条消息已被确认送达，再次发送便不会给地面中心带来新信息。本文的第二项和第三项研究分别对应图中的内容选择与发送交付环节。

在判断效果的评价指标上，分类准确率表示地面中心预测结果与参考类别一致的任务比例。然而，在五分类任务中不同类别的样本数量并不均等，如果模型只偏向预测样本数较多的类别，总体准确率仍可能偏高，但对少数类别的识别能力很差。为此，本文同时引入平衡准确率与宏平均 F1，其中前者计算各类别召回率的均值，后者为各类别精确率与召回率调和平均的算术均值，以便更全面地衡量各类别的识别表现。此外，语义选择研究将平均信念质量与平均送达概率的乘积定义为有效信念质量，用来同时反映语义判断水平和无线链路传输质量；交付控制研究则直接采用更严格的按时正确完成率，要求地面中心在截止时间前实际收到非空证据并作出正确判断。

在通信开销指标上，本文主要统计每任务平均词元（token）数、交付时延以及通信能耗。其中，词元数是文本经分词处理后的基本长度单位，用来衡量语义消息的文本开销，而实际无线链路传输的字节数还会受编码方式与协议报头影响；交付时延包含数据发送和等待确认的总耗时；单位正确任务通信能耗定义为总通信耗电与按时正确完成任务数之比。比较不同方法的节能效果时，均需同时结合按时正确完成率来看，避免把因放弃发送导致的耗电减少误认为性能提升。

\section{任务目标与通信约束}
\label{sec:base:task-objective}

设第 $k$ 项任务经内容筛选确定的候选消息集为 $\mathcal M_k$，传输结束时地面中心实际成功收到的不重复消息集合为 $S_k\subseteq\mathcal M_k$。地面中心依据 $S_k$ 给出最终预测 $\hat y_k(S_k)$，对应的参考类别为 $y_k^*$，任务实际完成时刻与截止时间分别为 $T_k$ 和 $D_k$。将单次任务按时正确完成定义为同时满足收到非空证据、判断正确且未超过时限三个条件，则连续任务流的按时正确完成率表示为
\begin{equation}
R_{\rm c}=\frac{1}{N_{\rm task}}\sum_{k=1}^{N_{\rm task}}
\mathbf1\{|S_k|>0,\ \hat y_k(S_k)=y_k^*,\ T_k\le D_k\}.
\label{eq:base:generic-control}
\end{equation}
式~\eqref{eq:base:generic-control} 将前端消息筛选与底层传输交付直接联系起来。相同数量的消息若包含的信息不同，可能导致截然不同的判断结果；而即便前端选出了关键消息，只要在无线传输中丢失或超时送达，该任务依然不能算作按时正确完成。

对于通信资源约束，记节点 $i$ 在整个任务序列中的累计通信耗电为 $E_i^{\rm comm}$，初始可用通信电量为 $E_i^{(0)}$；记消息 $m$ 的长度为 $\ell_m$，第 $k$ 项任务允许的消息总长度预算为 $B_k$，则系统运行须满足
\begin{equation}
E_i^{\rm comm}\le E_i^{(0)},\qquad
\sum_{m\in\mathcal M_k}\ell_m\le B_k.
\label{eq:base:resource-constraints}
\end{equation}
其中，长度预算主要限制前端候选筛选的总规模，实际发送与重传过程则进一步受节点剩余电量和任务剩余时间制约。提高发射功率虽然可以缩短单次发送耗时，但会加快节点耗电；多选候选或反复重传虽能增加判断把握，也会挤占后续任务的可用电量。因此，任何通信开销是否划算，最终都要结合式~\eqref{eq:base:generic-control} 的按时正确完成率来综合评估。

结合上述变量，本文三项研究的分工和参数联系十分明确。底层路由与功率控制决定链路速率与单次发送能耗，中层语义消息选择生成候选集合 $\mathcal M_k$ 及其长度，传输层发送与重传控制依据确认反馈安排发送时机，最终决定地面中心实际收到的集合 $S_k$ 和完成时间 $T_k$。表~\ref{tab:base:research-map} 汇总了三项研究在系统中的主要作用与相互依赖关系。

\begin{table}[H]
\centering\small
\bicaption{三个研究问题在观测信息回传中的作用与联系}{Roles and information dependencies of the three research problems}
\label{tab:base:research-map}
\begin{tabularx}{\linewidth}{@{}p{18mm}YYY@{}}
\toprule
比较维度&路由与功率控制&语义消息选择&发送与重传控制\\
\midrule
核心问题&持续回传中的队列积压与中继耗电&有限预算下的内容重叠与遮挡干扰&无线丢包与确认延迟或丢失\\
主要输入&邻居状态、队列积压、链路质量与剩余电量&候选语义、中心判断、链路代价与长度预算&候选消息、发送记录、累计确认与剩余资源\\
决策变量&下一跳节点及对应发射功率&消息来源、详略程度及停止时机&发送新消息、重传、等待或停止\\
模块衔接&提供链路速率与能耗模型，承接上层转发负载&输出筛选后的候选内容、来源及消息长度&沿用候选集与链路模型，决定最终接收集合\\
任务贡献&保障有限电量下的持续回传能力&优先挑选对正确判断最有帮助的信息&使关键证据在任务时限内可靠送达\\
评价重点&聚合吞吐量、丢包率与网络生命周期&分类准确率与平均通信词元长度&按时正确完成率与单位正确任务能耗\\
\bottomrule
\end{tabularx}
\end{table}

对应各自控制的决策变量，后文三章分别选取了直接反映该环节效果的评价指标。其中，吞吐量、丢包率与网络生命周期衡量连续回传的底层支撑能力，分类准确率衡量前端所选语义内容的判断价值，按时正确完成率与单位正确任务能耗则综合检验最终的交付结果和节能表现。同时，研究三直接以研究二筛选出的候选消息作为待发输入，并沿用研究一的链路速率与发送能耗计算方式，确保了不同章节在数据流和代价核算上保持一致。

\section{通信决策的求解基础}

本文的三项研究分别面临长期耗电约束、多条消息组合筛选以及接收状态无法直接观测三类难点。下面简要介绍后续章节所采用的李雅普诺夫优化、强化学习与接收状态估计方法。

当未来的数据到达和信道变化不可预知时，李雅普诺夫优化能够把跨时间的长期平均约束转化为每个时隙的本地在线决策\cite{neely2010stochastic}。对于长期平均能耗约束 $\limsup_{T\rightarrow\infty}T^{-1}\sum_{t<T}\E[e_i(t)]\le\bar e_i$，可以引入虚拟能量队列
\begin{equation}
Z_i(t+1)=\left[Z_i(t)+e_i(t)-\bar e_i\right]^+.
\label{eq:base:virtual-queue}
\end{equation}
虚拟队列 $Z_i(t)$ 用来记录截至当前时刻超出平均耗电额度的累计超支量。利用非负截断性质 $Z_i(t+1)\ge Z_i(t)+e_i(t)-\bar e_i$，对 $t=0,\ldots,T-1$ 累加可得
\begin{equation}
\frac{1}{T}\sum_{t=0}^{T-1}e_i(t)
\le \bar e_i+\frac{Z_i(T)-Z_i(0)}{T}.
\label{eq:base:finite-budget}
\end{equation}
只要虚拟队列随时间增长慢于线性速度，右侧余项在长时域下便会趋近于零，从而保证长期平均耗电不超过设定额度。在有限任务周期内，终端项 $Z_i(T)/T$ 不会完全消失，因此实际系统中虚拟队列主要用来平滑耗电节奏，而节点电量不透支的硬边界仍由剩余电量非负条件严格保证。

进一步定义二次李雅普诺夫函数 $L(\bm Z)=\frac12\sum_iZ_i^2$。在单步耗电有界的前提下，展开条件漂移 $\Delta(t)=\E[L(\bm Z(t+1))-L(\bm Z(t))\mid\bm Z(t)]$，可以得到包含交叉项 $\sum_iZ_i(t)(e_i(t)-\bar e_i)$ 的上界。将该漂移减去权重为 $V$ 的瞬时服务收益，即构成漂移加惩罚目标函数，用于在当前传输收益与长期资源透支之间动态权衡。第~\ref{rq1:chapter} 章将据此设计分布式路由与功率联合控制算法。

在语义消息选择中，多条消息融合后的判别效果受遮挡、视角重叠及先后顺序影响，难以用简单的数学公式直接表达，因此适合借助强化学习从交互反馈中学习序贯选择策略\cite{sutton2018rl}。与通过估计动作价值再间接选动作的深度 Q 网络（Deep Q-Network，DQN）及双重深度 Q 网络（Double DQN）不同\cite{mnih2015dqn,vanhasselt2016double}，本文采用策略梯度方法直接输出各候选动作的选择概率。决策状态 $s_h$ 汇总了候选特征、地面中心当前判断与剩余预算，并随每次选择动态更新；奖励函数则同时计入判断质量的提升和通信开销。为了防止重复发送同一消息或超出长度预算，引入可行动作集 $\mathcal A(s)$ 对策略网络输出的对数几率 $z_a$ 施加动作掩码，将非法动作的选取概率直接置零
\begin{equation}
\pi_\theta(a\mid s)=
\frac{\exp(z_a)\ind\{a\in\mathcal A(s)\}}
{\sum_{a'\in\mathcal A(s)}\exp(z_{a'})}.
\label{eq:base:masked-policy}
\end{equation}
其中停止动作始终位于可行集内，保证在没有可用候选或继续发送不再划算时能够正常结束选择。通过掩码拦截非法动作并借助奖励函数权衡合法候选，策略能够在满足硬约束的前提下优化整体回报。

策略参数采用近端策略优化（Proximal Policy Optimization，PPO）算法训练，通过对新旧策略概率比进行裁剪来限制单步更新幅度，避免训练过程剧烈波动\cite{ppo}
\begin{equation}
L^{\rm clip}(\theta)=\E\left[
\min\left(\rho_h(\theta)\widehat A_h,
\operatorname{clip}(\rho_h(\theta),1-\epsilon,1+\epsilon)\widehat A_h\right)\right],
\label{eq:base:ppo}
\end{equation}
其中 $\rho_h$ 表示新旧策略的概率比，$\widehat A_h$ 为优势估计，$\epsilon$ 为裁剪系数。训练中同时配合价值网络估计状态价值，并采用广义优势估计（Generalized Advantage Estimation，GAE）平衡估计偏差与方差\cite{gae}。第~\ref{rq2:chapter} 章将以此为基础，结合候选间的语义关联与已选历史设计具体的策略网络。

在发送与重传控制中，由于确认消息可能延迟到达或丢失，发送端收不到确认并不意味着数据没有送达，仅凭最新到达的确认信息无法确切知道接收端的真实状态。为此，借鉴部分可观测马尔可夫决策过程（Partially Observable Markov Decision Process，POMDP）的思路，用条件概率分布概括全部历史发送与反馈信息\cite{kaelbling1998pomdp}。记地面中心真实的接收状态为 $x_h$，截至步 $h$ 的发送与确认历史为 $\mathcal H_h$，则各可能接收状态的后验概率称为接收信念
\begin{equation}
b^{\rm rec}_h(x)=\Prb\{x_h=x\mid\mathcal H_h\}.
\label{eq:base:reception-belief}
\end{equation}
已知状态转移概率 $P(x'\mid x,a)$ 时，若当前未收到新确认，发送端按 $\bar b(x')=\sum_xP(x'\mid x,a)b(x)$ 向前递推更新接收信念；一旦收到新的累计确认，则先将确认生成时刻的历史接收状态固定下来，再依据此后发出的未确认记录重新推算当前各状态的概率分布。

这里的接收信念与前文的语义信念含义不同。语义信念表示地面中心对目标属于各个类别的判断概率，而接收信念表示发送端对地面中心已收到哪些消息的估计。得到当前接收状态的概率分布后，发送端还需据此决定后续动作。如果将未来多步中所有可能的发送、等待以及丢包和确认情况全部展开，需要评估的状态分支会随步数呈指数级增长，在线计算开销过高。为此，本文借鉴模型预测控制的做法\cite{mayne2000mpc}，每次只向前推算有限几步内发送与等待的效果，选出当前最合适的一步动作执行，待收到新的确认反馈后再更新接收状态并重新计算。

\section{本章小结}

本章建立了多无人机协同观测信息回传的系统通信模型，给出了通信能量、消息长度与任务时限三类约束以及任务按时正确完成的判定条件，并说明了路由与功率控制、语义消息选择和发送重传控制三个环节的分工与参数联系。最后，本章简要介绍了后续章节所用的李雅普诺夫优化、带动作掩码的策略优化以及接收状态估计方法。下一章将首先研究动态拓扑和有限电量下的无人机路由与功率联合控制方法。